\section*{Chapitre \ref{ch:b1:extent-q-and-k} --- Décrire un système chimique~: avancement, activités, Q et K}

\begin{solution}{exo:b1:extent-q-and-k:1}
\omterm{def:b1:extent-q-and-k:variables}{Intensives}~: température, \omterm{def:b1:extent-q-and-k:composition}{concentration molaire}, masse volumique,
pression, \omterm{def:b1:extent-q-and-k:composition}{fraction molaire}. \omterm{def:b1:extent-q-and-k:variables}{Extensives}~: volume, quantité de matière,
masse.
\end{solution}

\begin{solution}{exo:b1:extent-q-and-k:2}
$p(\ce{N2}) = 0.7808 \times 1.013 = \qty{0.791}{bar}$, $p(\ce{O2}) =
0.2095 \times 1.013 = \qty{0.212}{bar}$, $p(\ce{Ar}) = 0.0093 \times 1.013
= \qty{0.0094}{bar}$.
\end{solution}

\begin{solution}{exo:b1:extent-q-and-k:3}
$n(\ce{H2}) = 3.0 - 2\xi$, $n(\ce{O2}) = 2.0 - \xi$, $n(\ce{H2O}) =
2\xi$. $\xi_{\max} = \min(3.0/2,\ 2.0/1) = \qty{1.5}{mol}$~: le
dihydrogène est limitant. État final~: \ce{H2} 0, \ce{O2}
\qty{0.5}{mol}, \ce{H2O}
\qty{3.0}{mol}.
\end{solution}

\begin{solution}{exo:b1:extent-q-and-k:4}
(a) $Q = p_{\ce{CO2}}/p^\circ$. (b) $Q = (c^\circ)^2/([\ce{Ag+}][\ce{Cl-}])$.
(c) $Q = (p_{\ce{NH3}}/p^\circ)^2/[(p_{\ce{N2}}/p^\circ)(p_{\ce{H2}}/p^\circ)^3]$.
(d) $Q = [\ce{NH4+}][\ce{OH-}]/([\ce{NH3}]\,c^\circ)$, l'eau ayant une
\omterm{def:b1:extent-q-and-k:activity}{activité} égale à 1.
\end{solution}

\begin{solution}{exo:b1:extent-q-and-k:5}
$\mathrm{A \rightleftharpoons C}$ est la somme~: $K = 10 \times 0.5 = 5$. $\mathrm{C \rightleftharpoons A}$~: $1/5 =
0.2$. $\mathrm{2\,A \rightleftharpoons 2\,C}$~: $5^2 = 25$.
\end{solution}

\begin{solution}{exo:b1:extent-q-and-k:6}
$Q = 5.0 \times 10^{-5} < K^\circ$~: sens direct. $Q = 0.10 > K^\circ$~:
sens inverse. Réactifs seuls~: $Q = 0 < K^\circ$, sens direct.
\end{solution}

\begin{solution}{exo:b1:extent-q-and-k:7}
$n(\ce{A2}) = 1 - \xi$, $n(\ce{A}) = 2\xi$, total $1 + \xi$. Avec $p =
p^\circ$, $Q = \dfrac{(2\xi/(1+\xi))^2}{(1-\xi)/(1+\xi)} = \dfrac{4\xi^2}{1 -
\xi^2}$. $Q = 0.25$ donne $4.25\,\xi^2 = 0.25$, $\xi = \qty{0.243}{mol}$.
\end{solution}

\begin{solution}{exo:b1:extent-q-and-k:8}
Avec $x = [\ce{C}]$~: $x/(0.10 - x)^2 = 100$, donc $100x^2 - 21x + 1 = 0$
et $x = (21 - \sqrt{41})/200 = \qty{0.0730}{mol/L}$ (l'autre racine
dépasse 0.10). $[\ce{A}] = [\ce{B}] = \qty{0.0270}{mol/L}$, $[\ce{C}] =
\qty{0.0730}{mol/L}$, $\tau = 0.73$.
\end{solution}

\begin{solution}{exo:b1:extent-q-and-k:9}
Avec $n$ de chacun~: $Q = \xi^2/(n - \xi)^2 = K^\circ$, donc
$\xi/(n - \xi) =
\sqrt{K^\circ}$ et $\tau = \xi/n = \sqrt{K^\circ}/(1 + \sqrt{K^\circ})$.
$\tau = 0.99$ exige $\sqrt{K^\circ} = 99$, $K^\circ \approx \num{9.8e3}$~:
environ $10^4$, le seuil habituel d'une \omterm{def:b1:extent-q-and-k:quantitative}{réaction quantitative}.
\end{solution}

\begin{solution}{exo:b1:extent-q-and-k:10}
À l'équilibre, $p_{\ce{CO2}} = 0.20\,p^\circ$, c'est-à-dire $n(\ce{CO2}) =
pV/RT = 0.20 \times 10^5 \times 0.0100/(8.314 \times 1100) =
\qty{0.0219}{mol}$. Avec \qty{0.010}{mol} de carbonate, cette valeur ne
peut être atteinte~: tout le carbonate se décompose, $p_{\ce{CO2}} = 0.010 \times 8.314 \times
1100/0.0100 = \qty{9.15e3}{Pa} = \qty{0.091}{bar}$, $Q < K^\circ$, et
l'état final (\qty{0.010}{mol} de \ce{CaO}, \qty{0.010}{mol} de
\ce{CO2}, pas de \ce{CaCO3}) n'est pas un équilibre. Avec
\qty{0.050}{mol}~: équilibre, \qty{0.0219}{mol} de \ce{CO2} et de
\ce{CaO}, \qty{0.0281}{mol} de \ce{CaCO3} restant.
\end{solution}

\begin{solution}{exo:b1:extent-q-and-k:11}
À l'équilibre, $n(\ce{B}) = 2\,n(\ce{A})$ et $n(\ce{C}) = 3\,n(\ce{A})$
(le volume se simplifie). Alors $6\,n(\ce{A}) = 1.0$~: $n(\ce{A}) =
\qty{0.167}{mol}$, $n(\ce{B}) = \qty{0.333}{mol}$, $n(\ce{C}) =
\qty{0.500}{mol}$.
\end{solution}

\begin{solution}{exo:b1:extent-q-and-k:12}
Quantités égales~: $\xi^2/(0.50 - \xi)^2 = 4$, $\xi/(0.50 - \xi) = 2$,
$\xi = \qty{0.333}{mol}$, \omterm{def:b1:extent-q-and-k:fractional-extent}{rendement} $0.333/0.50 = 67\,\%$. Avec
\qty{1.0}{mol} d'alcool~: $\xi^2 = 4(0.50 - \xi)(1.0 - \xi)$,
$3\xi^2 - 6\xi + 2 = 0$,
$\xi = (6 - \sqrt{12})/6 = \qty{0.423}{mol}$, \omterm{def:b1:extent-q-and-k:fractional-extent}{rendement} 85\,\%~: un
excès d'un réactif augmente le \omterm{def:b1:extent-q-and-k:fractional-extent}{rendement} calculé sur l'autre.
\end{solution}

\begin{solution}{pb:b1:extent-q-and-k:1}
\textbf{1.} $x(\ce{CH4}) = 1.00/4.00 = 0.250$~; $x(\ce{H2O}) = 0.750$.
\textbf{2.} \qty{7.5}{bar} et \qty{22.5}{bar}.
\textbf{3.} Masses 16.04 et \qty{54.06}{g}~: $w(\ce{CH4}) = 0.229$,
$w(\ce{H2O}) = 0.771$.
\textbf{4.} Les \omterm{def:b1:extent-q-and-k:composition}{fractions molaires}, les \omterm{def:b1:extent-q-and-k:composition}{fractions massiques}, les
\omterm{def:b1:extent-q-and-k:composition}{pressions partielles} (et la pression totale)~; les quantités de matière
sont \omterm{def:b1:extent-q-and-k:variables}{extensives}.
\textbf{5.} $a = p_i/p^\circ$~: 7.5 pour le méthane, 22.5 pour la vapeur
d'eau.
\textbf{6.} \ce{CH4} $1.00 - \xi$, \ce{H2O} $3.00 - \xi$, \ce{CO} $\xi$,
\ce{H2} $3\xi$, total $4.00 + 2\xi$.
\textbf{7.} $\xi_{\max} = \qty{1.00}{mol}$~; le méthane est limitant.
\textbf{8.} \ce{CH4} 0, \ce{H2O} \qty{2.00}{mol}, \ce{CO} \qty{1.00}{mol},
\ce{H2} \qty{3.00}{mol}.
\textbf{9.} \qty{6.00}{mol}, contre \qty{4.00}{mol} dans l'alimentation~:
la réaction forme quatre molécules à partir de deux.
\textbf{10.} \ce{H2O} 0.333, \ce{CO} 0.167, \ce{H2} 0.500.
\textbf{11.} Pour que tout le méthane réagisse et pour laisser de la
vapeur d'eau à la réaction de conversion, dont elle abaisse le quotient.
\textbf{12.} \ce{CO} $1 - \xi$, \ce{H2O} $2 - \xi$, \ce{CO2} $\xi$, \ce{H2}
$3 + \xi$, total 6.00 quel que soit l'\omterm{def:b1:extent-q-and-k:extent}{avancement}.
\textbf{13.} $Q = \dfrac{(p_{\ce{CO2}}/p^\circ)(p_{\ce{H2}}/p^\circ)}
{(p_{\ce{CO}}/p^\circ)(p_{\ce{H2O}}/p^\circ)}$ avec $p_i = (n_i/6)p$~:
les facteurs $p/(6p^\circ)$ se simplifient entre numérateur et
dénominateur (deux molécules de gaz de chaque côté), ce qui laisse
l'expression donnée.
\textbf{14.} $\xi = 0$~: $Q = 0 < 4.0$, sens direct.
\textbf{15.} $\xi(3 + \xi) = 4(1 - \xi)(2 - \xi)$, c'est-à-dire
$3\xi^2 - 15\xi +
8 = 0$~: $\xi = (15 - \sqrt{129})/6 = \qty{0.607}{mol}$~; l'autre
racine, 4.39, dépasse $\xi_{\max} = 1$.
\textbf{16.} \ce{CO} \qty{0.393}{mol}, \ce{H2O} \qty{1.393}{mol},
\ce{CO2} \qty{0.607}{mol}, \ce{H2} \qty{3.607}{mol}.
\textbf{17.} $\tau = 0.607/1.00 = 0.607$.
\textbf{18.} $\xi(3 + \xi) = 4(1 - \xi)(5 - \xi)$, $3\xi^2 - 27\xi + 20 =
0$, $\xi = \qty{0.814}{mol}$~: plus de vapeur, plus de conversion.
\textbf{19.} $K^\circ = (0.99 \times 3.99)/(0.01 \times 1.01) = 391$.
\textbf{20.} $Q(\xi)$ est croissante, donc une constante $K^\circ$ plus
grande est atteinte pour un \omterm{def:b1:extent-q-and-k:extent}{avancement} plus grand. La constante de cette
réaction est plus grande à basse température, là où la réaction est
lente~: un premier réacteur travaille chaud et vite, un second plus
froid achève la conversion (l'influence de la température sur $K^\circ$
est traitée dans le volume de deuxième année).
\textbf{21.} $3.607/(6.00 - 1.393) = 3.607/4.607 = 0.783$.
\textbf{22.} \qty{0.393}{mol} de \ce{CO} et \qty{0.607}{mol} de \ce{CO2}.
\textbf{23.} $x(\ce{H2}) = 3.607/6.00 = \textbf{0.60}$~: six molécules sur
dix qui sortent du réacteur de conversion sont du dihydrogène.
\end{solution}
